\chapter{Неоднородные условия на торце и асимптотика решения}

\section{Постановка задачи}

Сохраняем индивидуальные типы условий: первый род при $x=0$,
второй род при $x=l$. Рассмотрим первую задачу в общем виде
\[
\begin{cases}
 u_t=a^2u_{xx}-bu+f(x,t),&0<x<l,\quad t>0,\\
 u(0,t)=g(t)=A\cos(\omega_1t)+B\sin(\omega_2t),\\
 -u_x(l,t)=0,\\
 u(x,0)=\phi(x).
\end{cases}
\]
Здесь $a>0$, $b\ge0$, $A,B\in\mathbb R$, $\omega_1,\omega_2>0$.
Поскольку правое условие однородно, оно эквивалентно $u_x(l,t)=0$.

В отличие от предыдущей задачи, первая задача допускает боковой
теплообмен ($b>0$) и распределённый источник $f$.
Поэтому ниже сначала приводится решение для произвольных $f,\phi$.
Исследование воздействия только граничной неоднородности выполняется
при $f=0$; иначе асимптотику определяет также источник.
Для теплоизолированной боковой поверхности дополнительно полагаем
$b=0$. Начальную температуру $\phi(x)=x$ можно подставить во все формулы.

\section{Замена и метод Фурье}

Положим
\[
 u(x,t)=g(t)+v(x,t).
\]
Тогда
\[
\begin{cases}
 v_t=a^2v_{xx}-bv+F(x,t),\\
 v(0,t)=0,\quad v_x(l,t)=0,\\
 v(x,0)=\phi(x)-g(0),
\end{cases}
\qquad F=f-g'-bg.
\]
Собственные функции и скорости затухания:
\[
 X_n(x)=\sin(k_nx),\quad k_n=\frac{(n+\tfrac12)\pi}{l},\quad
 \sigma_n=a^2k_n^2+b>0,\qquad n=0,1,\ldots.
\]
Ищем $v=\sum_{n=0}^{\infty}v_n(t)X_n(x)$.
Обозначим
\[
 d_n=\frac2l\int_0^lX_n(x)\,dx=\frac{2}{lk_n},\qquad
 f_n(t)=\frac2l\int_0^lf(\xi,t)\sin(k_n\xi)\,d\xi,
\]
\[
 c_n=\frac2l\int_0^l\bigl(\phi(\xi)-g(0)\bigr)\sin(k_n\xi)\,d\xi.
\]
Получаем
\[
 v_n'+\sigma_nv_n=f_n-d_n(g'+bg),\qquad v_n(0)=c_n.
\]
Следовательно, решение исходной задачи равно
\[
\boxed{\begin{aligned}
 u(x,t)=g(t)+\sum_{n=0}^{\infty}\sin(k_nx)
 \bigg[&c_ne^{-\sigma_nt}\\
 &+\int_0^te^{-\sigma_n(t-\tau)}
 \bigl(f_n(\tau)-d_n(g'(\tau)+bg(\tau))\bigr)\,d\tau\bigg].
\end{aligned}}
\]
Для заданной гармонической неоднородности
\[
 g'+bg=A\bigl(b\cos\omega_1t-\omega_1\sin\omega_1t\bigr)
 +B\bigl(b\sin\omega_2t+\omega_2\cos\omega_2t\bigr).
\]
Если $\phi(x)=x$, то $g(0)=A$ и
\[
 c_n=\frac2l\left(\frac{(-1)^n}{k_n^2}-\frac A{k_n}\right).
\]
Несогласованность начальных данных с граничными значениями в углах
не препятствует решению при $t>0$; начальное условие можно понимать
в смысле $L^2(0,l)$.

\section{Гармонический установившийся режим}

Пусть $f=0$. Введём комплексную функцию
\[
 q_\omega=\frac{\sqrt{b+i\omega}}{a},\qquad
 H_\omega(x)=\frac{\cosh(q_\omega(l-x))}{\cosh(q_\omega l)}.
\]
Берётся ветвь корня с положительной действительной частью.
Функция удовлетворяет условиям
\[
 a^2H_\omega''-(b+i\omega)H_\omega=0,\qquad
 H_\omega(0)=1,\quad H_\omega'(l)=0.
\]
Поэтому точное частное решение имеет вид
\[
\boxed{P(x,t)=A\operatorname{Re}\bigl(H_{\omega_1}(x)e^{i\omega_1t}\bigr)
 +B\operatorname{Im}\bigl(H_{\omega_2}(x)e^{i\omega_2t}\bigr).}
\]
Оно удовлетворяет уравнению, условию $P(0,t)=g(t)$ и $P_x(l,t)=0$.
Полное решение:
\[
\boxed{u(x,t)=P(x,t)+\sum_{n=0}^{\infty}r_n
 e^{-\sigma_nt}\sin(k_nx),}
\]
\[
 r_n=\frac2l\int_0^l\bigl(\phi(\xi)-P(\xi,0)\bigr)
 \sin(k_n\xi)\,d\xi.
\]
Таким образом,
\[
 \|u(\cdot,t)-P(\cdot,t)\|_{L^2}
 \le e^{-\sigma_0t}\|\phi-P(\cdot,0)\|_{L^2},\qquad
 \sigma_0=b+\frac{a^2\pi^2}{4l^2}.
\]
После любого положительного времени сглаживание также даёт
равномерное затухание переходной части.

При ненулевых амплитудах и положительных частотах решение в общем
случае не имеет предела: оно приближается к ограниченному колебательному
режиму. Если активна одна частота или отношение двух активных частот
рационально, режим периодический. При двух ненулевых амплитудах и
иррациональном отношении частот он квазипериодический.
Начальные условия влияют только на затухающую часть.

Следовательно, при фиксированной форме $A\cos\omega_1t+B\sin\omega_2t$
и $f=0$ неограниченный режим невозможен. При положительных частотах
ненулевой постоянный предел также невозможен. Выбор $A=B=0$
даёт затухание к нулю. Если допускается $\omega_1=0$, косинусный член
становится постоянным и возникает стационарный профиль, рассмотренный ниже.
Чтобы получить все требуемые типы поведения, далее подбираем другие
временные зависимости $g(t)$, сохраняя типы границ и правое условие нулевым.

\section{Аналитический подбор воздействий}

Во всех трёх примерах этой части считаем $f=0$, а начальную функцию
$\phi$ произвольной, в частности $\phi(x)=x$.

\subsection{Решение стремится к нулю}

Выберем ненулевое затухающее воздействие
\[
 g(t)=C e^{-rt},\qquad C\ne0,\qquad 0<r<\sigma_0.
\]
Определим $H_r$ из
\[
 a^2H_r''-(b-r)H_r=0,\qquad H_r(0)=1,\quad H_r'(l)=0.
\]
Явно,
\[
 H_r(x)=\begin{cases}
 \displaystyle\frac{\cosh(\kappa(l-x))}{\cosh(\kappa l)},
 &b>r,\quad\kappa=\sqrt{b-r}/a,\\[0.7em]
 1,&b=r,\\[0.4em]
 \displaystyle\frac{\cos(\eta(l-x))}{\cos(\eta l)},
 &b<r,\quad\eta=\sqrt{r-b}/a.
 \end{cases}
\]
Условие $r<\sigma_0$ исключает нуль знаменателя.
Точное частное решение $P_r=C e^{-rt}H_r$ даёт
\[
 u(x,t)=C e^{-rt}H_r(x)+\sum_{n=0}^{\infty}z_n
 e^{-\sigma_nt}\sin(k_nx),
\]
где $z_n$ --- коэффициенты Фурье функции $\phi-CH_r$.
Обе части затухают, поэтому
\[
\boxed{u(x,t)\longrightarrow0.}
\]

\subsection{Решение неограниченно возрастает}

Выберем
\[
 g(t)=C e^{rt},\qquad C>0,\quad r>0.
\]
Пусть
\[
 \kappa_r=\frac{\sqrt{b+r}}{a},\qquad
 K_r(x)=\frac{\cosh(\kappa_r(l-x))}{\cosh(\kappa_rl)}>0.
\]
Тогда $P_r=C e^{rt}K_r$ --- точное частное решение, и
\[
 u(x,t)=C e^{rt}K_r(x)+\sum_{n=0}^{\infty}z_n
 e^{-\sigma_nt}\sin(k_nx),
\]
где $z_n$ --- коэффициенты функции $\phi-CK_r$.
Поскольку $K_r$ имеет положительный минимум на $[0,l]$,
\[
\boxed{u(x,t)\longrightarrow+\infty}
\]
в каждой точке и равномерно по $x$. Неограниченность следует не только
из роста температуры на границе, но и из явного решения внутри стержня.

\subsection{Постоянный предельный режим}

Выберем $g(t)=C$, $C\ne0$. При $b=0$ точное частное решение равно $C$:
\[
 u(x,t)=C+\sum_{n=0}^{\infty}z_n
 e^{-a^2k_n^2t}\sin(k_nx),\qquad
 z_n=\frac2l\int_0^l(\phi(\xi)-C)\sin(k_n\xi)\,d\xi.
\]
Таким образом, для теплоизолированной боковой поверхности
\[
\boxed{u(x,t)\longrightarrow C.}
\]

При $b>0$ и $f=0$ то же граничное воздействие приводит не к
пространственной константе, а к стационарному профилю
\[
\boxed{U_C(x)=C\frac{\cosh(\sqrt b(l-x)/a)}{\cosh(\sqrt bl/a)}.}
\]
Переходная часть вновь является суммой затухающих мод.
Важно различать постоянство во времени и постоянство по координате.
Если требуется именно $U\equiv C\ne0$, то при $b>0$, $f=0$
такой профиль невозможен при любом выборе граничного воздействия:
подстановка в стационарное уравнение даёт $0=-bC$.
Для его поддержания нужен источник $f=bC$ и граничные данные
$u(0,t)=C$, $u_x(l,t)=0$. При этих изменённых данных $u-C$
решает однородную затухающую задачу.

\section{Роль распределённого источника}

При произвольном $f(x,t)$ поведение нельзя определить только по $g$.
К решению с $f=0$ добавляется
\[
 u_f(x,t)=\sum_{n=0}^{\infty}\sin(k_nx)
 \int_0^te^{-\sigma_n(t-\tau)}f_n(\tau)\,d\tau.
\]
В частности, постоянный источник обычно создаёт ненулевой стационарный
вклад, даже если граничное воздействие затухает; растущий источник
может создавать неограниченный режим при нулевых границах.
Поэтому приведённый аналитический подбор изолирует действие граничных
неоднородностей условием $f=0$ и явно различает случаи $b=0$ и $b>0$.
